\section{Plataformas, herramientas y la mente del usuario}

Una vez que el Agent puede hacer cosas, falta responder qué es en la mente del usuario. Esta pregunta define los límites del producto: si es una herramienta, una plataforma o un nuevo punto de entrada.

Zhang Yueguang (张月光) también habló de las oportunidades para nuevas plataformas. Según él, una nueva plataforma de contenido necesita un nuevo medio y una nueva forma de interacción. Sora no necesariamente se convertirá en plataforma, porque podría seguir siendo un formato de contenido de video corto. ChatGPT, en cambio, sí podría convertirse en plataforma, porque creó un nuevo punto de entrada para la interacción. En el caso de los Agent, términos como vertical/general o modelo/aplicación son demasiado gruesos: lo que realmente importa es con qué idea te recuerda el usuario.

\begin{figure}[H]
\centering
\includegraphics[width=0.94\textwidth]{platform-tool-agent-positioning.png}
\caption{Oportunidades para nuevas plataformas: productos de tipo herramienta, productos de tipo plataforma y el Agent como punto de entrada.}
\end{figure}

\begin{knowledgebox}{Lectura de la figura: la mente del usuario define los límites}
Una herramienta resuelve tareas concretas; una plataforma crea relaciones y un ecosistema. Un Agent como punto de entrada debe lograr que el usuario asocie la marca con una necesidad que tiene con frecuencia. Si el alcance es demasiado estrecho, se usa poco; si es demasiado amplio, «poder hacer de todo equivale a no poder hacer nada». Dokie se posiciona como un Agent de creación y expresión, no como una mezcla de todas las tareas posibles.
\end{knowledgebox}

\subsection{Los límites de la creación y la expresión}

Zhang Yueguang considera que las presentaciones (PPT), los documentos, la expresión y la creación pueden reunirse en una idea de nivel superior en la mente del usuario: ayudarlo a crear y a expresarse mejor. En cambio, tareas como recopilar información, llenar formularios o procesar datos no necesariamente se integrarán con un Agent de creación. Es un criterio de producto: los límites no se trazan según las capacidades del modelo, sino según la mente del usuario y los escenarios de uso.

\begin{importantbox}{El problema del posicionamiento de un Agent}
Un Agent no es mejor por ser más general. El producto debe encontrar en la mente del usuario un límite que este quiera recordar, que se use con frecuencia y que sea lo bastante amplio sin volverse vago.
\end{importantbox}

\subsection{Resumen de la sección}

La competencia entre productos de Agent no es solo una competencia de capacidades de los modelos, sino también por un lugar en la mente del usuario. Que pase de ser una herramienta a ser un punto de entrada depende de que tenga un límite de uso frecuente, una relación con el usuario y capacidad de seguir ampliándose.
